\section{Deification of Man}

In \textbf{Guido de Giorgio}, we often find the best from \textbf{Rene Guenon} and \textbf{Julius Evola}. Yet, the ``personal equation'' is qualitatively different. Guenon is cold and logical, the path of the jnani whose goal is the Truth. Evola sometimes writes with a deep passion, but is often wordy; his is the path of action, with the Good as its goal. De Giorgio, on the other hand, seems to speak as a man of conviction, arising from an inner knowing, while avoiding the sentimental and the merely personal. He is a poet and mystic, following the highest path, that of Absolute Beauty\footnote{\url{https://gornahoor.net/?p=3537}}.

Guenon and Evola, for all their virtues, are ultimately escapists, Guenon to the ``traditional'' East and Evola to the long forgotten Ancient Rome. But De Giorgio offers a different way. Unlike Guenon who seems to regard the various Traditions as discrete and independent, Evola and De Giorgio see a continuity in the Traditions of the West, from the Hyperboreans to the Medieval era. Yet Evola wants to stop part way, and even to reverse it, while De Giorgio wants to continue and complete the rectification that Evola described.

Taking Guenon's claim seriously that \textbf{Dante}'s Divine Comedy is a blueprint of the path for initiates, he describes a vision for the reconstitution of Tradition in his La Tradizione Romana, under the aspects of the Divine Cycle, the roles of the castes, the role of the Leader, the spirit of the Roman tradition, and its fulfillment in the Comedy.

In the following passage from Aforismi e Poesie (Archè Milano, 1999), De Giorgio describes the process of Deification, drawing on Traditional notions. From Guenon, we see that ``man'' is just one of the states of being, but does not exhaust all the states. From Evola, we see the theme of transcending man as such, and the way to know God is to be God\footnote{\url{https://gornahoor.net/?p=1875}}. But we suspect, however, that De Giorgio is saying the same thing, not as an abstract theory, but rather from his own interiority. Readers can judge for themselves. De Giorgio's text follows:

\begin{quotationx}
To lose consciousness of oneself, to become conscious of God. To become conscious of the nothing that I am in order to have consciousness of everything that is God. One does not take a step toward God without this preliminary negative. What I am not, God is.

Man makes himself God in the measure in which he is no longer man, that is, when his transparency, in the absolute sense, is such that God passes through him. This is the deification of man, not human imperfection arrogantly divinized---a haughtiness absolutely outside the place in life that leads to God---but man annulled, abolished, dissipated, so that God alone is because God alone is God and no god is god except God.
\end{quotationx}



\flrightit{Posted on 2012-05-03 by Cologero}

\begin{center}* * *\end{center}

\begin{footnotesize}
\begin{sffamily}
\texttt{Farid Othman al Deenari on 2012-05-04 at 02:50 said:}

``Guenon and Evola, for all their virtues, are ultimately escapists... ''

We object to this characterization. In fact, it is the world around us that is `escapist', whilst we, in such a milieu, are automatically faulted for being anchored in the realm of knowledge. This realm, being ever present and not subject to temporal conditions, namely space and time, does not allow one to escape from it; nor could we escape to it, for it is all that exists, both theoretically, and practically to those who possess effective consciousness of it.

Insofar as the West deviates, the defender of truth can and will be misrepresented as an `escapist' by simply holding his ground.

\hfill

\texttt{Cologero on 2012-05-04 at 05:23 said:}

Farid, I did not intend ``escapist'' as a personal insult to those men, but rather in the sense that, on a personal level, they abandoned (i.e., ``escaped'' from) the West. Whatever you may think about not being subject to temporal conditions, as a stark matter of fact, Guenon did indeed translate himself in Space, and Evola was nostalgic for a different Time.

Those of us who choose to hold our ground --- in the West at this historical moment --- cannot limit ourselves to their choices.

\hfill

\texttt{Angolmois on 2012-05-04 at 13:53 said:}

De Giorgio quote reminded me about the initiatic experience of being a complete human personality with all traits intact while witnessing also at the same time the universal consciousness.

Astrologically speaking this would mean the dual experience of Leo (integrated personality) and Aquarius (universal consciousness). De Giorgio seems to be taking a stance for the universal consciousness (`pratyeka'), while Evola was more of an alchemist and a hermetist in this regard, since he ``remained in the cosmos and got back to it'' instead of choosing to escape from it through the Capricornian gate.

Cologero: I find it hard to see either Guénon and Evola as escapists, yet I understand what you mean with it. And after all, did not Evola describe the way you mention in your last post in his `Ride the Tiger' with unfathomable precision? ``A way out, further in'' etc.?

\hfill

\texttt{Charlotte on 2012-05-04 at 19:09 said:}

woah, rewind, Capricorn gate, say more bro -- I've been there!! It's in Atlantis, it's scary, at least I was scared... .panic stations, password, aaarrggg!!

\hfill

\texttt{Charlotte on 2012-05-04 at 19:58 said:}

deep at the bottom of de-nial!!

\hfill

\texttt{Angolmois on 2012-05-05 at 04:39 said:}

Concerning Atlantis and escapism, did not the Masters of Light leave the degenerated Atlantis when they saw that it is too late and futile to save it wholly; the preservation of the core of the primordial tradition clearly outweighed the saving of a few souls who do not even recognise the need to be saved.

Charlotte: I've been in the temples Atlantis also in my dreams. Finely arranged beautiful temples and amazing surroundings, that's all I can say. The water was in a very important role in the surrounding gardens, and the whole scene was very beautiful. Although this is my personal fantasy, it certainly reflects the stories told about the Atlantean tradition.

\hfill

\texttt{Charlotte on 2012-05-05 at 06:44 said:}

well that's kind of right, but an essential factor is missing there. The Atlanteans, as you probably know, were of one mind -- not of `scattered' consciousness such as humanity experienced after the destruction of the Tower of Babel. This enabled them to do extraordinary things in terms of building, for instance, and the manifestation of `gods' in what is perhaps best understood in terms of the literary myths we now have. They were also very `in tune' (understatement) with the forces of nature, which is what led ultimately to their downfall. The power they had access to became corruptive within their souls and they began performing ever more extreme feats of power to do with the weather, for instance, and fire -- the discovery of nuclear energy, in particular. They also became very interested in sex and started to enhance their bodies in order to maximise sexual pleasure (hence their association now with the dolphin race, which is the most sexual creature on earth, humans included). This is why there are such uncanny resonances with the modern world, which is staring into its mirror at Atlantis -- a magical future prophecy as much as a past `legend', or buried memory (for `weather modification', read `HAARP'). As far as I'm aware the memory of the flood and all that went with it was so traumatic that humanity collectively wiped it out of its consciousness, leading to the state of `slumber' experienced by the mass of humankind. And yes of course it was the `elite', as ever, who largely escaped the consequences of the group actions, of which they -- as the ones in power -- were must culpable. Another parallel with modern life in which the `elites' -- all that means now is those with money -- are more able than most to escape the horrific consequences of `the fall'... .to a point... .there's only so far one -- even a privileged one -- can go along that road, be sure of it. There's more than one reason why the princes Buddha and Jesus renounced their wealth. The escape of the elites is a plan that has once again been activated, but this time around more people are aware of the situation -- more people whose memories of Atlantis are NOT buried. The priests weren't quite the only ones who escaped. They planned their exodus, but others simply survived and were washed up on other shores. A lot of these have ended up as `new agers' because of their resentment at the way the priestcraft handled the catastrophe -- causing the problem then saving themselves. So even though they `know' the hierarchy holds truths about God and man, there is a reaction against the hierarchy and a drive to make known what was formerly kept secret, so that everyone has a chance to survive. A lot of those who imagine themselvs to be elite now are not really `in the club'. They have been given a certain amount of information and this in itself is so monumental that for the most part they are content to simply indulge their password potential, opening of the magical doors and access to the other world. In this way are some of the greatest minds and initiates kept a degree lower than they are really able to achieve, because if they were to attain to the highest orders they would then understand the terrible side of the truth and comprehend the extent to which they have been fooled. The key at that point is to break open the closed circle by invoking the name of Jesus Christ and through love and compassion (Happy Birthday Buddha!)throwing one's lot in to help the children in the nursery to survive, to awaken, in other words. Tangentially, the swastika goes back to at least Atlantean times. It is still possible to go there in the astral so your dreams, if lucid, are most likely to be true. I've pestered Cologero about this so many times I see no reason to break the habit and hold back now -- here is my story :-)

\url{http://alchemical-weddings.com/alchemical-weddings/atlantis}


\url{http://alchemical-weddings.com/alchemical-weddings/half-man-half-fish}


\url{http://alchemical-weddings.com/alchemical-weddings/swastika}


\url{http://alchemical-weddings.com/alchemical-weddings/ancient-script}

\url{http://alchemical-weddings.com/alchemical-weddings/pisces-capricorn}


\hfill

\texttt{Charlotte on 2012-05-05 at 06:55 said:}

this isn't just `altruism', by the way. Humanity is bound together by its collective karma, so it's essential that the `weight' of this is light as a feather, that the balance is tipped towards ascension not sinking. United we'll stand and divided we fall... .

\hfill

\texttt{mleow on 2012-05-05 at 09:52 said:}

I also find the notion that ``Evola was nostalgic for a different Time'' is misleading. At least that is not the sense and understanding i get of the man and his writing, as Angolmois said ``Ride the tiger'' is perhaps the best example that Evola was in fact the one from the ``traditionalist school'' with the least escapist/nostalgic notions trying to cope with the modern reality as it is and all its destructive tendencies as incitments to inner awakening, instead of trying to conjure up artifical paradises, medieval, roman or whatever. The word ``Good'' in this regard also seems a bit strange, he always mentions the ``Unconditioned'' as the highest goal to steer ones compass to. Would you care to elaborate or refer to parts of his writing as regarding claims for this goal towards the ``Good''?

That said, im hoping to see more of these Giorgio translations. In what form has his writings survived is it merely texts from various correspondences or has he also written books on esoteric topics?

\hfill

\texttt{Cologero on 2012-05-05 at 10:28 said:}

The object of Action is the Good, otherwise action would be purposeless. Please think of it in the Traditional sense as a transcendental. For Plato, the Good itself was the ultimate. But as transcendentals, the Good, the True, and the Beautiful are ultimately unconditioned and One. How else could Action ``lead'' to the Unconditioned?

How is ``riding the tiger'' not passive? You are basing your judgment of Evola only on what has been made available to you. We have undertaken the effort to show the Evola of the 30s and 40s to demonstrate that he was decidedly not ``riding the tiger'', but engaged in the political situation of the time. A man may have to ride the tiger, when historico-material conditions are unsuited to him. But our position here is that conditions have changed, so to ride the tiger is no longer an option. \emph{A fortiori}, the time is come for the elite to arise and create the proper conditions. It is your option to \emph{cope} with conditions, and continue to whine about the modern world, but the point is, as someone once said, to change it.

De Giorgio published nothing while he was alive. \textit{La Tradizione Romana} existed in manuscript form in the 30s, presumably a copy was sent to Mussolini. It is systematic and describes the lineaments of the constitution of a future Traditional society. At that time he was associated with Evola and presumably they were trying to influence the political situation in the early 1930s toward a more Traditional direction.

At some point, De Giorgio withdrew, living a solitary life in the mountainous part of Italy, where he traveled a more mystical path. He left behind notebooks of poems and short essays, but nothing systematic; these have been gradually published after his death. I will translate more of the Roman Tradition on Gornahoor to show that Tradition is more than nostalgia. I will translate the more mystical and devotional texts elsewhere, since they may not be suited to the more analytical tone of Gornahoor. For example:

\textit{Woman}: \url{http://www.meditationsonthetarot.com/woman}, from \textit{Aphorisms and Poems}.

\hfill

\texttt{Charlotte on 2012-05-05 at 11:08 said:}

So what course will the `elite' take in trying to improve conditions in the present world for one and all, because a political party or philosophical school won't achieve this? What is required is positive practical action, of that there can be no doubt. Who is prepared to give up home comforts, abandon the supermarket in favour of small stores, cut out food additives, quit supporting big businesses, to cease buying into the system, to desist from believing the lies of politicians, to understand that the corporate tapeworm economy is strangling the world -- who can and WILL break free of the matrix-net surrounding us, be they elite or otherwise? Acknowledging with searing honesty who the desposynic heirs really are -- and what it really takes to partake of the promised spiritual truths -- is a difficult but necessary first step to achieving the desirable goal of saving the world. A leap of faith is required by almost everyone, plus individual determination to reform the self as a microcosmic alchemical procedure that brings about the cosmic renewal. Really there is no point denying it!

\hfill

\texttt{Angolmois on 2012-05-05 at 15:06 said:}

I have personally understood Evola's `Ride the Tiger' precisely as a path of action and positive change in-side the modern world. In Evola's own words it is meant for ``the elite'', but only as a handbook and a manual for individuals who may perhaps have no apparent connexion between each other. If this is out of the question for you Cologero and ``not enough'', I think that Gornahoor and websites / blogs like this one are one of the best way for the eventual organisation and appearance of the spiritual hiearchy of the earth that will change the world ``in an instant'' when the time is ready. Right now ``we'' are still scattered around and largely un-organised, not the least because of petty infighting among those who should first and foremost have a mutual understanding and the spirit of co-operation in their hearts.

\hfill

\texttt{mleow on 2012-05-06 at 04:11 said:}

That makes it more clear.

I have read all of Evola's work that are available in english and the few snippets of translations made avaible on Gornahoor has not made any drastic revelation or revaluation about him or his work on my part. Well please remember that ``ride the tiger'' was one of his later books so its quite natural that he did not pursued that path of apoletia etc under the war and fascist years, where at least the possibilty and ``humanmaterial'' existed to change and influence the millieu in a traditional direction. I dont whine much really anymore about being born in this declining world. I accept my situation and embrace it as Evola would say, as a spiritual test, try to see things according to truth and don't succumb or waste my life trying to cut of hydraheads ie spend my time on any vain action, as a spiritual master once said ``Who has improved the world? People have improved themselves and then gone.''

You say that the conditions have changed and yet in your next sentence you say that an elite will arise to create the new conditions? Sounds paradoxical to me. From my point of view i don´t see it as realistic that a spiritual elite will spring up and take control, could you please specify as concretely as possibly and not merely in the abstract how the current conditions seems apt for such an event to occur? How will this legion of beermongering footballcrazed pursuit of pleasurist sentimentalist secular cynical rationalist people ever recognise ,not to mention obey or need such an elite??

\hfill

\texttt{Charlotte on 2012-05-06 at 04:34 said:}

Let's keep it real shall we? the actual spiritual elite are out there right now saving lives, with no time to sit around chit-chatting, hampered mainly by the refusal of the super-rich and super-powerful to either offer support to the solution or cease working against God and for themselves only

\hfill

\texttt{Cologero on 2012-05-06 at 12:56 said:}

Evola's subtitle for Ride the Tiger was this: ``orientamenti esistenziali per un'epoca di dissoluzione''. Apparently, it was the American publishers who added: A Survival Manual for the Aristocrats of the Soul. Does anyone suppose that aristocrats of the soul publish ``how to'' manuals in mass market paperback editions? As the real Marcuse wrote, everything is co-opted by the system.

No, mleow, I cannot describe in detail, because it involves a singularity, the end of one world and the beginning of another. However, \emph{contra factum non est disputandum}, and the facts point to traditional civilizations, or new spiritual movements, arising seemingly out of nothing. But I agree with you about avoiding vain action, or disputing with the ignorant.

Charlotte, you do true religion a disservice by trying to restrict it to moralism and humanism. That has never been the case.

\hfill

\texttt{Charlotte on 2012-05-06 at 15:09 said:}

Since when did keeping it real mean doing `true religion' a disservice -- maybe the truth hurts? I thought works were the ultimate proof of faith -- the royal road is one we all love to tread, but the greatest Masters preach and practice poverty -- how can this be denied? I'd be the last one to discount the power of prayer but I'm not seeing much of that either and I don't like hypocrisy -- talk about `elites' makes me cringe and like it or not, it IS the humble practioners who tend to be the most touched by God's grace, you can see it shining from them. Not that I'm counting myself amongst this saintly group, just pointing out basic facts -- we are looong past the theoretical stage of forming the `new world', Utopia was good enough. Now it's time for implementation and in my opinion it's beyond urgent, but people are wholly resistant to change for the better if it means disrupting their own comfortable lives even by one iota. And I'm not a humanist, you know I am not, there is no need to be bitchy -- something about not shooting the messenger? Apologies, however, if I've over-moralising, I know that can be very annoying... .chit chat is fine, I'll butt out :-)

\hfill

\texttt{Cologero on 2012-05-06 at 18:51 said:}

I sense a severe communications breakdown here, but I'm not sure what it is. You speak of ``Masters'', who are, I believe, what we mean by \emph{spiritual elite}. Didn't Tomberg say somewhere that at the necessary moment, someone will appear to right things? I believe he mentioned by name St. Francis and St. Ignatius. Do you cringe because we single them out as ``saints''?

We have no objection to spiritual and corporal works, but please keep in mind the story of Martha and Mary. The present danger is a man-centred religion, when it needs to be God-centred. The contemplative life has always been primary and everything else flows therefrom.

\hfill

\texttt{Angolmois on 2012-05-07 at 03:17 said:}

I believe there is no need to despair if things seem to be currently out of hands concerning making things right again. The world will dissolve and break apart, that's the very promie of all scriptures, and then the ``elite'' or ``the masters'' -- who are preparing themselves as we speak -- will come out of the hermetic closet and reveal the real ``new world order'' and the new age ``in the spirit of Lif and Lifhtrasil''; life, light and breath of the spirit(ual hierarchy).

Having battled with this issue for years in desperation and in the dark night of the soul I simply some day realised that it actually requires quite little in reality for things to change for the better again. It requires the change in the hearts and minds of men, and for this we all must aim our prayers and deeds. The Good One is still in charge no matter how it looks like right now.

Blessings of The High One for you all.

\hfill

\texttt{Charlotte on 2012-05-07 at 05:10 said:}

Well I think it's the focus on Rome I'm finding difficult, but I accept that this is a large part of what Gornahoor is about so fair enough because it is your website. I certainly never said mention of `saints' made me cringe, the opposite in fact, I said mention of `elite', I don't like the connotations of that word, it reminds me of the nazis. I accept, however, that it may not be meant in the way I'm reading it. I believe true `Masters', like St Francis, practice poverty and love humanity, the notion of a superstructure of `more brilliant' minds and people controlling the ignorant masses is what `elitism' speaks of. I also think I make a valid point about works -- the `Masters' of the ilk of St Francis literally went out and healed people, as did the ultimate Master, Jesus. Perhaps I sense my own `lack', for want of a better word, in the face of people who are able to put into practice more of what they know, or who know `less' but do more to help humanity. Put it this way, a starving child needs prayers for sure, but also something else, right? Maybe I'm just reaching the stage of my spiritual journey where I'd like to put the food into their mouths more often, rather than `just' praying. Which again, is not to discount the power of prayer, more collective positive thinking in the world could work miracles, but in the meantime there are multiple crises that need tackling on multiple fronts. Another example of a `humble' soul of the kind I was referring to might be America's own George Washington Carver -- having read something of his life for the first time quite recently he certainly emerges as a `Master' initiate of sorts, so in this respect truly `elite'. Perhaps the real issue, then, is one of defining precisely what constitutes this `elite'. Maybe I've been misunderstanding by thinking that people here are singled out as being `elite' simply for having adopted a certain political/religious philosophy and going on to write about it? And yes, I do believe there is a Messiah who rights things, ultimately there is only one Master, the great leveller... .I agree the change must come in the hearts and midns of all humans, this is what I'm meaning by `grass roots' level action, whereby all individuals take responsibility. The time of teaching has been going on for a long while, it's time to put these things into practice. As for `New World Order'... .well, we had a recent discussion about the Georgia Guidestones. Unfortunately most of those who might think of themselves as `elite' in this context are simply rich and powerful. Any esoteric knowledge they might have gleaned has come from decidedly suspect angles, OTO, etc -- and all mixed in with the influx of nazism into America following WWII. They never really lost the war, you know?

\hfill

\texttt{Charlotte on 2012-05-07 at 05:13 said:}

I also refer to Einstin: `those who know have a duty to act'

\hfill

\texttt{Cologero on 2012-05-07 at 08:19 said:}

Perhaps, Charlotte, you could read \textit{The Second Precept of Love}\footnote{\url{http://www.meditationsonthetarot.com/the-second-precept-of-love}}, particularly what he says about the prehension of the All or Unity. We tend to experience the world in discrete pieces rather than in its totality. This applies to Time as well as Space; in other words, the panorama of History is also the unfolding of a divine Plan, not just the work of men.

Guenon complained about the change in the word ``elite'' from the way he first used it. It may be better to think of it as the ``elect'', as in ``many are called, few are chosen (electus)''. So, it is not something a person may claim for himself, but rather something that is assigned to him from above. Your Master also speaks of the leaven in the bread, so we are entirely faithful to that message. The change must come from the elect, long before it comes ``in the hearts and minds of all humans''.

Who are you, Charlotte, Martha or Mary?

\hfill

\texttt{Cologero on 2012-05-07 at 08:36 said:}

As for the focus on Rome, that is due to its central place in the history of Europe, a fortiori, the very creation and self-understanding of Europe, not once, but twice. Aeneas, the spiritual founder of Rome, brought the Wisdom of the East to the West, of which Rome was the center. Dante, in De Monarchia, points out the Divine Plan for Ancient Rome, and I assume everyone is familiar with that.

Then we see the central place of Rome in the Tarot, in the Arcana of the High Priest of Rome and the Emperor of Rome. As their influence wanes, the world falls apart.

By the way, there are no Nazis in USA, quite the contrary. USA is the heir to Cromwell's England and the British Freemasons. OTO is a British, not a Nazi, organization.

\hfill

\texttt{Charlotte on 2012-05-07 at 09:04 said:}

Yes I know, from Crowley, but a lot of the `magick' employed is similar to the way the nazis achieved their dark spiritual aims. A lot of nazi scientists, as you know, escaped to America, they were taken on board by the government. When I say `they won', I mean in terms of the technological/material emphasis in the world today, in particualr the warcraft. Hitler took heart to some extent -- for his racial theories, at least -- from the American wars against the Indian, and he was also fairly admiring of the English, it's all tied up together as I'm sure you're aware. What is helpful, however, is a clarification of what is meant by `elite', thank you -- as this is a term that's often bandied about these days in a very negative sense.This is an open website so who knows what casual readers might think? (To some extent I'm one of those, if also regular). I still maintain, however, that the change has already come from the minds of the `elect', the forerunners, and it is up to us now to implement this in some way. To bring it out of the world of formation and into the world in a way that is recognisable to all. The way of salvation. Rome, well, I'm `Greek', as I think you also know -- my tutor was very snooty about Roman history in comparison and I agree with him almost entirely on this point! The answer is to revive the Mysteries, drink from the lake of memory -- here!
\url{http://alchemical-weddings.com/alchemical-weddings/lake-of-memory-2}

\hfill

\texttt{Angolmois on 2012-05-07 at 10:52 said:}

For the sake of clarification, Charlotte \& others, what I meant by `the new world order' has certainly nothing to do with what is meant by it today. I used it simply to refer to the emerging Tradition in the ruins of the current world, a transcendental order and a new golden age, the age of Saturn.

The Elite defines itself and \textit{shines}, its activity is salvific in all orders of life according to the dharma of the particular practicioner or order. The crown and the sceptre, on the other hand, always move only further away from those who seek and yearn for them.

\hfill

\end{sffamily}
\end{footnotesize}

